% Research proposal for a grant or funding call.
% Edit the metadata below, then replace the section text. Layout, fonts and
% the accent colour live in style.tex. Build with pdflatex and biber.
\documentclass[11pt]{article}

% --- Metadata: change these for your proposal --------------------------------
\newcommand{\ProjectTitle}{Community-Grounded Speech and Text Technology for Southern Quechua and Aymara}
\newcommand{\ShortTitle}{Qillqa}
\newcommand{\ProposalCall}{Language Futures Programme 2027}
\newcommand{\PI}{Dr Inés Condori Huamán}
\newcommand{\Institution}{Andes Institute for Language and Computation, Cusco}
\newcommand{\Duration}{36 months, from 1 April 2027}
\newcommand{\Amount}{EUR 1,060,000}

% Layout and typography for the proposal. main.tex holds the content;
% the accent colour is the one setting most people change here.

% --- Type: Source Serif body, Source Sans display, matching math ------------
\usepackage[T1]{fontenc}
\usepackage[osf]{sourceserif}
\usepackage[lining,tabular,semibold,scale=0.94]{sourcesans}
\usepackage{amsmath}
\usepackage{newtxmath}
\usepackage[italic,defaultmathsizes,noendash]{mathastext}
\usepackage[protrusion=true,expansion=true,tracking=true,kerning=true]{microtype}
\SetTracking{encoding=*,shape=*,family=*}{0}

% --- Page: A4 with a measure of roughly 70 characters ------------------------
\usepackage[a4paper,textwidth=146mm,top=32mm,bottom=30mm,headsep=10mm,
  footskip=14mm,hcentering]{geometry}
\linespread{1.08}
\setlength{\parindent}{1.4em}
\setlength{\parskip}{0pt}
\frenchspacing
\clubpenalty=10000
\widowpenalty=10000

% --- Colour: change the accent here -----------------------------------------
\usepackage{xcolor}
\definecolor{accent}{HTML}{1D5A63}   % deep teal
\colorlet{accentsoft}{accent!7}
\colorlet{accentmid}{accent!45}
\colorlet{muted}{black!58}
\colorlet{hairline}{black!22}

% --- Small labels: tracked capitals in the sans -----------------------------
\newcommand{\caps}[1]{{\sffamily\textls[110]{\MakeUppercase{#1}}}}

% --- Headings ---------------------------------------------------------------
\usepackage[explicit]{titlesec}
\titleformat{\section}
  {\sffamily\large}
  {\llap{\color{accent}\thesection\hspace{0.9em}}}
  {0pt}{#1}
\titlespacing*{\section}{0pt}{2.4\baselineskip plus 4pt minus 2pt}{0.7\baselineskip}
\titleformat{\subsection}
  {\sffamily\normalsize\color{accent}}
  {\thesubsection\hspace{0.7em}}
  {0pt}{#1}
\titlespacing*{\subsection}{0pt}{1.4\baselineskip plus 2pt minus 2pt}{0.35\baselineskip}

% --- Running header and footer ----------------------------------------------
\usepackage{lastpage}
\usepackage{fancyhdr}
\pagestyle{fancy}
\fancyhf{}
\renewcommand{\headrulewidth}{0pt}
\fancyhead[L]{\footnotesize\color{muted}\caps{\ShortTitle}}
\fancyhead[R]{\footnotesize\color{muted}\sffamily\ProposalCall}
\fancyfoot[C]{\footnotesize\color{muted}\sffamily
  Page \thepage\ of \pageref*{LastPage}}
\fancypagestyle{cover}{\fancyhf{}\fancyfoot[C]{\footnotesize\color{muted}%
  \sffamily Page \thepage\ of \pageref*{LastPage}}}

% --- Lists, tables, captions ------------------------------------------------
\usepackage{enumitem}
\setlist{itemsep=0.2\baselineskip,topsep=0.4\baselineskip,leftmargin=1.6em}
\usepackage{booktabs,tabularx,array,float}
\usepackage[group-separator={,},group-minimum-digits=4,
  text-family-to-math=true]{siunitx}
\newcolumntype{L}{>{\raggedright\arraybackslash}X}
\setlength{\heavyrulewidth}{0.7pt}
\setlength{\lightrulewidth}{0.35pt}
\setlength{\cmidrulewidth}{0.35pt}
\renewcommand{\arraystretch}{1.25}
\newcommand{\tablefont}{\sffamily\small}
\newcommand{\thead}[1]{\leavevmode{\color{accent}\footnotesize\caps{#1}}}
\usepackage[font={small,sf},labelfont={color=accent},labelsep=quad,
  justification=raggedright,singlelinecheck=false,skip=8pt]{caption}

% --- Panels: summary and aims -----------------------------------------------
\usepackage[most]{tcolorbox}
\newtcolorbox{summarybox}{enhanced,breakable,colback=accentsoft,
  colframe=accentsoft,boxrule=0pt,arc=0pt,left=6mm,right=6mm,top=4.5mm,
  bottom=4.5mm,before skip=6mm,after skip=4mm,
  overlay={\node[anchor=north west,inner sep=0pt]
    at ([xshift=6mm,yshift=-4.5mm]frame.north west)
    {\footnotesize\color{accent}\caps{Summary}};},
  fontupper=\small\linespread{1.12}\selectfont,
  before upper={\vspace*{1.6em}\noindent}}

\newcounter{aim}
\newcommand{\aimobjective}{\item}
\NewDocumentEnvironment{aim}{m}
  {\refstepcounter{aim}%
   \begin{tcolorbox}[enhanced,breakable,colback=white,colframe=white,
     boxrule=0pt,arc=0pt,left=5mm,right=0mm,top=1mm,bottom=1mm,
     before skip=4mm,after skip=4mm,
     borderline west={1.2pt}{0pt}{accent}]
   {\footnotesize\color{accent}\caps{Aim \theaim}}\par\nobreak
   \vspace{0.2\baselineskip}%
   {\sffamily\large #1\par}\nobreak\vspace{0.35\baselineskip}}
  {\end{tcolorbox}}
\newlist{objectives}{enumerate}{1}
\setlist[objectives]{label={\sffamily\small\color{accent}O\theaim.\arabic*},
  ref={O\theaim.\arabic*},leftmargin=3em,labelsep=0.8em,
  itemsep=0.15\baselineskip,topsep=0.3\baselineskip}

% --- Gantt chart ------------------------------------------------------------
\usepackage{tikz}
\usetikzlibrary{shapes.geometric}

% --- Bibliography and cross references --------------------------------------
\usepackage[backend=biber,style=authoryear-comp,maxcitenames=2,
  maxbibnames=6,uniquelist=false,dashed=false]{biblatex}
\renewcommand{\bibfont}{\small}
\setlength{\bibitemsep}{0.35\baselineskip}
\setlength{\bibhang}{1.4em}
\DeclareFieldFormat[article,inproceedings,incollection]{title}{#1}
\renewbibmacro{in:}{\ifentrytype{article}{}{\printtext{\bibstring{in}\intitlepunct}}}
\addbibresource{refs.bib}

\usepackage[colorlinks,linkcolor=accent,citecolor=accent,urlcolor=accent,
  pdfusetitle]{hyperref}
\usepackage[capitalise,noabbrev]{cleveref}
\crefname{aim}{Aim}{Aims}
\crefname{objectivesi}{Objective}{Objectives}
\crefname{wp}{work package}{work packages}
\urlstyle{same}


\title{\ShortTitle: \ProjectTitle}
\author{\PI}

\begin{document}
\thispagestyle{cover}

% --- Cover block --------------------------------------------------------------
\begingroup
\noindent{\footnotesize\color{accent}\caps{Research proposal \quad\textperiodcentered\quad \ProposalCall}}\par
\vspace{9mm}
\noindent{\fontsize{30}{32}\selectfont\color{accent}\ShortTitle}\par
\vspace{3mm}
\noindent\begin{minipage}{0.86\linewidth}\raggedright
  \fontsize{17}{22}\selectfont\ProjectTitle
\end{minipage}\par
\vspace{9mm}
\noindent{\color{hairline}\rule{\linewidth}{0.4pt}}\par
\vspace{3mm}
\noindent\begin{tabularx}{\linewidth}{@{}>{\raggedright\arraybackslash}p{0.38\linewidth}
    >{\raggedright\arraybackslash}X >{\raggedright\arraybackslash}p{0.22\linewidth}@{}}
  {\scriptsize\color{muted}\caps{Principal investigator}} &
  {\scriptsize\color{muted}\caps{Duration}} &
  {\scriptsize\color{muted}\caps{Amount requested}} \\[-1pt]
  \PI & \Duration & \Amount \\[5pt]
  \multicolumn{3}{@{}l}{\scriptsize\color{muted}\caps{Host institution}} \\[-1pt]
  \multicolumn{3}{@{}l}{\Institution} \\
\end{tabularx}\par
\vspace{2mm}
\noindent{\color{hairline}\rule{\linewidth}{0.4pt}}\par
\endgroup

\begin{summarybox}
Southern Quechua and Aymara are spoken by more than nine million people, yet
neither has a speech recogniser, spell checker or machine translation system
that its speakers would choose to use. \ShortTitle{} will build these tools
together with three speaker communities in Cusco, Puno and La Paz. Over three
years we will record and transcribe 600 hours of consented speech, release an
open text corpus under community data governance, adapt multilingual models
to both languages with less than a tenth of the data such models usually
require, and evaluate every system with speakers rather than with automatic
scores alone. The project ends with tools that run offline on an ordinary
phone, and with a governance model that other indigenous language
communities can adopt.
\end{summarybox}

\section{Background}\label{sec:background}

Language technology now reaches most of the world's speakers, but not most
of the world's languages. Of the roughly 7,000 languages in use, fewer than a
hundred have the text and speech resources on which current systems are
trained \parencite{ortega2024resource}. Quechua and Aymara are large by any
demographic measure \parencite{adelaar2004languages}, yet they sit firmly on the wrong side of this divide:
the largest public Quechua text collection holds under four million words,
much of it translated religious material \parencite{huaman2023corpus}.

Two recent developments make progress possible. First, massively multilingual
models transfer surprisingly well to languages they have barely seen, provided
the target language shares script and some vocabulary with the training
mixture \parencite{lindqvist2025transfer}. Second, community-led data
initiatives have shown that speakers will contribute recordings in large
numbers when they keep control over how the data is used
\parencite{tupa2022governance,carroll2020care}. What is missing is a project
that joins the two: one that treats speakers as partners in design and
evaluation, not only as a source of data.

Evaluation is the weakest link. Automatic metrics such as word error rate say
little about whether a transcription is usable, and they are unreliable for
agglutinative languages, where a single morpheme error changes a whole word
\parencite{mendoza2024metrics}. \ShortTitle{} therefore makes human evaluation
by speakers a first-class part of the work plan (\cref{sec:workplan}).

\section{Aims}\label{sec:aims}

The project pursues three aims. Each is tied to one or more work packages in
\cref{fig:gantt}, and each has objectives that a reviewer can check at the end
of the grant.

\begin{aim}{Build open, consented resources}\label{aim:data}
  Create speech and text collections for both languages under a data
  governance framework that the partner communities write and enforce.
  \begin{objectives}
    \item Sign data agreements with three community organisations before any
          recording starts.\label{obj:agreements}
    \item Record and transcribe 600 hours of speech across six dialect areas.
    \item Release a 25-million-word text corpus with per-document provenance.
  \end{objectives}
\end{aim}

\begin{aim}{Adapt models with little data}\label{aim:models}
  Show that multilingual models can reach usable quality in Quechua and Aymara
  with a fraction of the data they normally need.
  \begin{objectives}
    \item Reach a word error rate below 18\,\% on conversational speech.
    \item Match a strong supervised baseline in translation to and from
          Spanish using under 200,000 sentence pairs.
    \item Distil every model to run offline on a phone with 4\,GB of memory.
  \end{objectives}
\end{aim}

\begin{aim}{Evaluate with speakers}\label{aim:evaluation}
  Replace automatic scores as the final judge with structured evaluation by
  speakers of each dialect, and publish the protocol.
  \begin{objectives}
    \item Run four rounds of speaker evaluation, each with at least 60
          participants.
    \item Publish an evaluation protocol that other projects can reuse.
  \end{objectives}
\end{aim}

\section{Methods}\label{sec:methods}

\subsection{Data collection and governance}

Recording follows the protocol of \textcite{tupa2022governance}: speakers
choose the topics, keep the right to withdraw their recordings, and decide
together which uses are permitted. Transcription is done by trained community
research assistants, paid at the regional rate for skilled work, with a
second assistant checking a random 10\,\% sample. No recording is collected
before \cref{obj:agreements} is met.

\subsection{Model adaptation}

We start from an open multilingual speech encoder and an open translation
model, and adapt both in two stages. The first stage continues pre-training on
untranscribed audio and monolingual text. The second fine-tunes on labelled
data with a loss that weights each dialect $d$ by the inverse of its share
$p_d$ of the training set, so that minority dialects are not drowned out:
\begin{equation}\label{eq:loss}
  \mathcal{L} = \sum_{d \in D} \frac{p_d^{-\alpha}}{\sum_{d'} p_{d'}^{-\alpha}}
  \; \mathbb{E}_{(x,y) \sim \mathcal{D}_d} \bigl[ -\log p_\theta(y \mid x) \bigr],
\end{equation}
where $\alpha \in [0,1]$ controls how strongly rare dialects are favoured
\parencite{arivazhagan2019massively}. We tune $\alpha$ on held-out speakers,
never on held-out sentences from speakers seen in training.

\subsection{Evaluation}

Each evaluation round pairs automatic metrics with a speaker study. Speakers
rate transcriptions and translations for adequacy and naturalness on a
five-point scale, and mark errors that would change the meaning. We report
agreement between raters and between ratings and \cref{eq:loss}-trained model
scores, following \textcite{mendoza2024metrics}.

\section{Work plan}\label{sec:workplan}

The work is organised in five work packages over twelve quarters
(\cref{fig:gantt}). Governance runs through the whole project, because data
agreements are revisited at every release.

\begin{figure}[t]
\centering
\begin{tikzpicture}[x=7.7mm,y=-6.6mm,font=\sffamily\footnotesize,
    bar/.style={fill=accent,draw=none},
    softbar/.style={fill=accentmid,draw=none},
    milestone/.style={diamond,draw=accent,fill=white,line width=0.8pt,
      inner sep=1.7pt}]
  % label column ends at x = 0; quarters run from 0 to 12
  \foreach \y/\yr in {1/{Year 1},2/{Year 2},3/{Year 3}}
    \node[text=accent] at ({4*\y-2},-1.05) {\caps{\yr}};
  \foreach \q in {1,...,12}{
    \pgfmathtruncatemacro{\qq}{mod(\q-1,4)+1}
    \node[text=muted,font=\sffamily\scriptsize] at (\q-0.5,-0.3) {Q\qq};
  }
  \foreach \x in {0,4,8,12}
    \draw[hairline,line width=0.4pt] (\x,-1.45) -- (\x,5.6);
  \foreach \x in {1,2,3,5,6,7,9,10,11}
    \draw[black!8,line width=0.3pt] (\x,0.15) -- (\x,5.6);
  \draw[black!60,line width=0.5pt] (-6.5,0.15) -- (12,0.15);
  \draw[black!60,line width=0.5pt] (-6.5,5.6) -- (12,5.6);

  \foreach \row/\wp/\name in {
      1/WP1/{Governance and partnership},
      2/WP2/{Speech and text collection},
      3/WP3/{Model adaptation},
      4/WP4/{Speaker evaluation},
      5/WP5/{Release and dissemination}}{
    \node[anchor=west,text=accent] at (-6.5,\row) {\wp};
    \node[anchor=west] at (-5.5,\row) {\name};
  }
  \fill[softbar] (0,0.8) rectangle (12,1.2);
  \fill[bar] (0,0.8) rectangle (2,1.2);
  \fill[bar] (1,1.8) rectangle (7,2.2);
  \fill[bar] (4,2.8) rectangle (10,3.2);
  \fill[softbar] (6,3.8) rectangle (12,4.2);
  \fill[bar] (6,3.8) rectangle (7,4.2);
  \fill[bar] (8,3.8) rectangle (9,4.2);
  \fill[bar] (10,3.8) rectangle (11,4.2);
  \fill[bar] (11.5,3.8) rectangle (12,4.2);
  \fill[bar] (8,4.8) rectangle (12,5.2);

  \foreach \x/\row/\m in {2/1/M1,6/2/M2,8/3/M3,11/4/M4,12/5/M5}{
    \node[milestone] at (\x,\row) {};
    \node[anchor=south,text=accent,font=\sffamily\scriptsize]
      at (\x,\row-0.26) {\m};
  }
\end{tikzpicture}
\caption{Work plan by quarter. Dark bars mark intensive work and pale bars
continuing work; in WP4 the dark blocks are the four speaker evaluation
rounds. Diamonds mark the milestones listed in \cref{tab:milestones}.}
\label{fig:gantt}
\end{figure}

\begin{table}[H]
\centering\tablefont
\caption{Milestones and the quarter in which each is due.}\label{tab:milestones}
\begin{tabularx}{\linewidth}{@{}l L r@{}}
  \toprule
  \thead{Milestone} & \thead{Deliverable} & \thead{Due} \\
  \midrule
  M1 & Data agreements signed with all three partner communities & Q2 \\
  M2 & First public release of the speech and text corpus & Q6 \\
  M3 & Adapted speech and translation models reach the targets of \cref{aim:models} & Q8 \\
  M4 & Fourth speaker evaluation round completed and analysed & Q11 \\
  M5 & Offline tools, corpus and evaluation protocol released & Q12 \\
  \bottomrule
\end{tabularx}
\end{table}

\section{Risks and mitigation}\label{sec:risks}

\Cref{tab:risks} lists the risks we consider most likely to delay the
project, with the step we take before each one occurs and the fallback if it
does.

\begin{table}[H]
\centering\tablefont
\caption{Main risks, their assessed likelihood and impact, and how each is
mitigated.}\label{tab:risks}
\begin{tabularx}{\linewidth}{@{}>{\raggedright\arraybackslash}p{0.27\linewidth}
    l l L@{}}
  \toprule
  \thead{Risk} & \thead{Likelihood} & \thead{Impact} & \thead{Mitigation} \\
  \midrule
  Data agreements take longer than planned & Medium & High &
    Negotiation starts six months before the grant; WP3 can begin on
    existing public text. \\
  Too few speakers in one dialect area & Medium & Medium &
    Recruit through radio stations and schools; merge with a neighbouring
    area after Q4 if needed. \\
  Models miss the error-rate target & Low & High &
    Report the gap openly; shift effort from translation to speech, where
    gains are largest. \\
  Loss of a community research assistant & Medium & Low &
    Train two assistants per site, so that each site keeps one trained
    member at all times. \\
  Compute costs rise during the project & Low & Medium &
    Budget includes a 15\,\% contingency on compute; distilled models
    reduce later training cost. \\
  \bottomrule
\end{tabularx}
\end{table}

\section{Budget}\label{sec:budget}

The amount requested is \Amount{} over three years (\cref{tab:budget}).
Personnel makes up 74\,\% of direct costs; most of the rest pays speakers
and community workshops, which we treat as part of the research, not as
overhead.

\begin{table}[H]
\centering\tablefont
\caption{Budget by year, in euros. Indirect costs are the programme's
flat rate of 25\,\% of direct costs.}\label{tab:budget}
\begin{tabular*}{\linewidth}{@{\extracolsep{\fill}}l
    S[table-format=7.0] S[table-format=7.0] S[table-format=7.0]
    S[table-format=7.0]@{}}
  \toprule
  \thead{Cost category} & {\thead{Year 1}} & {\thead{Year 2}} &
    {\thead{Year 3}} & {\thead{Total}} \\
  \midrule
  Postdoctoral researchers (2.0 FTE)   & 118400 & 121900 & 125600 & 365900 \\
  Doctoral researcher (1.0 FTE)        &  38200 &  39300 &  40500 & 118000 \\
  Community research assistants        &  46800 &  48200 &  49600 & 144600 \\
  Speaker compensation and workshops   &  22500 &  28000 &  18500 &  69000 \\
  Recording equipment                  &  14800 &   2000 &      0 &  16800 \\
  Travel and fieldwork                 &  18600 &  21400 &  16200 &  56200 \\
  Compute and storage                  &   9500 &  24000 &  21500 &  55000 \\
  Dissemination and open access        &   2000 &   6500 &  14000 &  22500 \\
  \cmidrule(l){2-5}
  Direct costs                         & 270800 & 291300 & 285900 & 848000 \\
  Indirect costs (25\,\%)              &  67700 &  72825 &  71475 & 212000 \\
  \midrule
  {\color{accent}Total requested}      & 338500 & 364125 & 357375 & 1060000 \\
  \bottomrule
\end{tabular*}
\end{table}

\section{Expected impact}\label{sec:impact}

For speakers, \ShortTitle{} delivers a speech recogniser, a spell checker and
a translation aid that work without a network connection, in the dialects
people actually speak. For the field, it provides the first evaluation
protocol for these languages that is grounded in speaker judgement, and
evidence on how far multilingual transfer can be pushed with community-scale
data \parencite{nakamura2026distil}. For other indigenous language
communities, the governance agreements and their record of use over three
years offer a tested model, in the spirit of the CARE principles
\parencite{carroll2020care}, that can be adopted without starting from
nothing. All software is released under an open licence; the data is
released under the terms each community sets in \cref{aim:data}.

\printbibliography[heading=bibnumbered,title={References}]

\end{document}
